\section{Observability and Monitoring}
\label{sec:observability}

Monitoring in this platform serves two distinct goals with different requirements. Operationally, it must convert failures from user-visible to administrator-visible: the incident analysed in Section~\ref{sec:driver-mismatch} was detected through a failed user spawn precisely because no node-level GPU telemetry existed, and the alerting introduced since is designed to close that gap. Methodologically, time-series data on GPU utilization, scheduling behavior, and service health constitutes the empirical foundation for evaluating the resource-efficiency and operability claims of Chapter~\ref{cha:system-architecture-design}; without continuous measurement, such claims are anecdotal. Monitoring is therefore treated as a first-class platform component rather than an operational afterthought, following the practice standardized in large-scale production systems, where the four golden signals of latency, traffic, errors, and saturation form the core of service health assessment~\cite{beyer2016sre}.

\subsection{Stack Selection and Deployment}

The metrics stack is built from the kube-prometheus-stack distribution of the Prometheus ecosystem. Prometheus~\citep{volz2015prometheus} descends from Google's Borgmon monitoring system used alongside the Borg cluster manager~\citep{burns2016borg}; it consolidates a multi-dimensional time series database (TSDB), pull-based scraping with service discovery, and a dedicated query language into a single open-source system, and has become the de facto metrics standard in the Kubernetes ecosystem. Three properties determined the choice. First, the Prometheus Operator extends the Kubernetes API with custom resources (\texttt{ServiceMonitor}, \texttt{PrometheusRule}), so that scrape targets and alerting rules are expressed declaratively and managed by the same Git-driven lifecycle as every other platform component. Second, the exporter ecosystem already covers the layers this platform must observe: \texttt{node-exporter} for host metrics, \texttt{kube-state-metrics} for the state of Kubernetes objects, DCGM for GPU telemetry, and native \texttt{/metrics} endpoints of JupyterHub, SeaweedFS, and Traefik. Third, the stack is sufficiently compact to fit the resource envelope of the service node, consistent with the resource-efficiency objective of Chapter~\ref{cha:system-architecture-design}.

Grafana fronts the stack for dashboards and is integrated into the Keycloak identity layer like every other user-facing service, so that access to operational data requires no separate credentials. Alertmanager receives firing alerts from Prometheus rule evaluation; alert delivery is currently limited to the dashboards themselves, which matches the single-administrator scale of the platform, and can be extended to email or chat webhooks without structural changes.

\subsection{Tiered Metric Storage}

The storage philosophy layers media by access pattern and failure semantics rather than using a single backend. The Prometheus write-ahead log (WAL) and TSDB head require a local filesystem with crash-safe writes: a network filesystem mounted through FUSE, such as the SeaweedFS CSI storage class used for notebook volumes, cannot guarantee WAL integrity across mount interruptions, and documented production incidents show WAL corruption and query failures when the metrics database itself resides on such a mount. The hot tier therefore runs on node-local storage with a 30-day retention window, which covers routine operational lookback.

Long-term retention instead uses the SeaweedFS S3 object interface, its strongest and most native access mode. A Thanos sidecar~\citep{thanos} ships compacted two-hour blocks from the local TSDB to a dedicated \texttt{monitoring} bucket, where they accumulate for the duration of the platform's lifetime; this is the data source for the utilization analyses of the evaluation, surviving local rotation and node rebuilds. The layering preserves a deliberate independence property: the monitoring system must not depend on the infrastructure it monitors. If SeaweedFS is unavailable, block uploads retry while the local buffer absorbs the delay, and no scrape path traverses the storage layer. The same pattern, local hot state with SeaweedFS S3 for durable bulk data, mirrors the platform-wide data lifecycle decisions for registry images and experiment artifacts, reusing one storage system through the interface it provides best.

\subsection{GPU Telemetry and Driver-Health Alerting}

GPU nodes run the NVIDIA DCGM exporter~\citep{dcgm}, which exposes per-GPU counters (utilization, memory occupancy, temperatures, ECC and XID error events) through the standard \texttt{/metrics} endpoint. Deployment requires the NVIDIA container runtime to inject the host driver libraries and device nodes into the exporter pod; without this injection the exporter starts and serves metrics with an empty result, a failure mode that is silent to liveness probes. The heterogeneous GPU fleet constrains the exporter version: current DCGM releases drop support for the Pascal architecture of the TITAN~Xp, so the deployed 3.3 series is retained to keep both GPU models fully observable.

The primary alert, \texttt{DCGMExporterDown}, fires on exporter liveness rather than on individual counters. This design targets precisely the incident class of Section~\ref{sec:driver-mismatch}: a broken host driver prevents NVML initialization, the exporter crash-loops, and the alert fires independently of whether any GPU workload is running, converting the outage from user-visible to administrator-visible. A complementary rule alerts on XID error events, which indicate hardware and driver-level faults. In the division of labor argued in Section~\ref{sec:agentic-ops}, these alerts act as the surfacing mechanism for both human and AI operators: they shorten the detection phase, which the incident analysis identified as the phase most improved by tooling rather than by further automation of actions.
